\documentclass[10pt]{article}

\usepackage[utf8]{inputenc}

\usepackage[T1]{fontenc}

\usepackage{amsmath}

\usepackage{amsfonts}

\usepackage{amssymb}

\usepackage[version=4]{mhchem}

\usepackage{stmaryrd}

\usepackage{bbold}



\begin{document}

тируемо. Известно, что если в теории $E^{*}$ ориентируемы все векторныө расслоения, то



\[

E^{*}(X) \approx H^{*}\left(X ; \tilde{E}\left(S^{0}\right)\right)

\]



причем $2 E^{*}\left(S^{0}\right)=0$. В этом контексте иногда ослабляют условия на теорию $E^{*}$, напр. снимают условие коммутативности умножения и т. Д. Для любой теории $E^{*}$, в к-рой все комплексные расслоения ориентируемы, имеется гомоморфизм теории $U^{*} \rightarrow E^{*}$, где $U^{*}-$ теория унитарных кобордизмов, и этот гомоморфизм полностью задается $E$-ориентацией канонич. расслоения $\eta$ над $\mathbb{C} P^{\infty}$. Аналогичное верно и для $S_{p}$-расслоений (см. Кобордизм). Построение для любого класса векторных расслоений универсальной теории, отображающейся в любую теорию, где ориентируем данный класс расслоений, еще (1983) не проведено.



О р и е т т а ц и е й (другое название - ф у н д ам е н т а льны й к л а с с) замкнутого n-мерного многообразия (или, более общо, комплекса Пуанкаре формальной размерности $n$ ) в теории $E^{*}$ наз. такой элемент $z \in E_{n}(M)$, что гомоморфизм $E^{i}(M) \rightarrow E_{n-i}(M)$ вида $x \rightarrow z \bigcap x$ (см. [9]) есть изоморфизм; это - т. н. изоморфизм двойственности Пуанкаре. Оказывается, что многообразие (комплекс Пуанкаре) $E$-ориентируемо тогда и только тогда, когда $E$-ориентируемо его нормальное расслоение. Определяется $\mathrm{O}$. и для многообразий (комплексов Пуанкаре) с краем.



Лum. : [1] Д у б ров ин Б. А., Но в иков С. П., Фоме нк о А. Т., Современная геометрия, М., 1979; [2] Введение в топологию, М., 1980; [3] Р о х л и н В. А., Ф у к с Д. Б., Начальный курс топологии. Геометрические главы, М., 1977; [4] Х Х о ом о л л е р Д., Расслоенные пространства, пер. с англ., М., 1970; [5] С п е н е р Э., Алгебраическая топология, пер. с англ., M., 1971; [6] Д о л в Д А., "Математика", 1965 , т. 9, № 2, c. 8-, 1977; [8] С т о н г Р. Заметки по теории кобордизмов пер., англ., М., 1973; [9] у а й т х е д Дж., Новейшие достижения в теории гомотопий, пер. с англ., М., 1974.



ОРИСФЕРА - поверхность пространства Лобачевского, ортогональная к прямым, параллельным в некотором направлении. О. можно рассматривать как сферу с бесконечно удаленным центром. На О. реализуется евклидова геометрия, если под прямыми понимать орициклы, порядок точек определить через порядок Іірямых в пучке параллелей, определяющих орицикл, а движением называть такие движения в пространстве Јобачевского, к-рые переводят О. в себя.



А. Б. Иванов.



ОРИЦИКЛ, п р е д е льн а я ли н и я,- ортогональная траектория параллельных в нек-ром направлении прямых плоскости Лобачевского. О. можно рассматривать как окружность с бесконечно удаленным центром. О., порожденные одним пучком параллельных Ірямых, конгруэнтны, концентричны (т. е. высекают на прямых пучка конгруәнтные отрезки), незамкнуты и вогнуты в сторону параллельности прямых пучка. Кривизна О. постоянна. В модели Пуанкаре О.- окружность, касающаяся изнутри абсолюта.



Прямая и О. либо не имеют общих точек, либо касаются, либо пересекаются в двух точках под равными углами, либо пересекаются в одной точке под прямым углом.



Через две точки плоскости Лобачевского гроходят два и только два $\mathrm{O}$



Лит.: [1] К а г а н В. Ф., Основания геометрии, ч. 1-2, М. - Л., 1949-56; [2] Н о р д ен А. П., Элементарное введение в геометрию Јобачевского, М., 1953 ; [3] Е Ф и м о в Н. В., Высшая геометрия, 6 изд., М., 1978. А. Б. Иванов.



ОРИЦИКЛИЧЕСКИЙ ПОТОК - поток в пространстве биэдров такого $n$-мерного риманова многообразия $M^{n}$ (обычно замкнутого), для к-рого определено понятие орицикла; О. п. описывает движение биэдров вдоль определяемых ими орициклов. Основные случаи, когда определено понятие орицикла,- те, когда кривизна римановой метрики отрицательна и либо $n=2$, либо кривизна постоянна. Б и ә д $\mathbf{p} \mathbf{y}$, т. е. ортонормированному 2 -реперу $\left(x, e_{1}, e_{2}\right)\left(x \in M^{n}\right.$; $e_{1}, e_{2}$ - взаимно ортогональные единичные касательные векторы в точке $x$ ), сопоставляется орицикл $h\left(x, e_{1}, e_{2}\right)$, к-рый проходит через $x$ в направлении $e_{2}$ и расположен на проходящей через $x$ орисфере $H\left(x, e_{1}\right)$, являющейся $(n-1)$-мерным ортогональным многообразием семейства геодезич. линий, асимптотических (в положительном направлении) к геодезич. линии, проходящей через $x$ в ваправлении $e_{1}$. Направление на $h$, определяемое $e_{2}$, принимается за положительное (при $n=2$ роль $e_{2}$ только к этому и сводится; $H$ и $h$ могут иметь самопересечения; простейший способ избежать могущих возникнуть из-за этого неясностей состоит в том, чтобы выполнить аналогичные построения не в $M^{n}$, а в его универсальном накрывающем многообразии - при постоянной кривизне это обычное $n$-мерное пространство Лобачевского - и спроектировать полученный там орицикл в $M^{n}$ ). Под действием О. п. биэдр $\left(x, e_{1}, e_{2}\right)$ за время $t$ переходит в



\[

\left(x(t), e_{1}(t), e_{2}(t)\right),

\]



гце $x(t)$ с изменением $t$ движется с единичной скоростью по $h\left(x, e_{1}, e_{2}\right)$ в положительном направлении, единичный вектор $e_{1}(t)$ 'ортогонален $H\left(x, e_{1}\right)$ в точке $x(t)$ (выбор одного из двух возможных направлений для $e_{1}(t)$ производится по непрерывности) и $e_{2}(t)=d x(t) / d t$.



Изучение О. п. было начато в связи с тем, что он играл важную роль при исследовании геодезических потоков на многообразиях отрицательной кривизны [1]. Позднее эта роль перешла к нек-рым слоения., возникаюпим в теории $\boldsymbol{V}$-систем, а О. п. стал самостоятельным объектом исследования. Свойства О. П. хорошо изучены (см. [2] - [7], [11]). О нек-рых обобщениях см. в [8] - [10].



Лum.: [1] X о п ф Э., “Успехи матем. наук», 1949, т. 4, в. 2, c. 129-70; [2] П а р а с юк О. С., там же, 1953, т. 8, в. 3, с. $125-26$; [3] Г у р е в и ч Б. М., "Докл. АН'СССР", 1961 , т. 136, № 4, c. 768-70; [4] F u r s t e n be r g H., B KH.: Recent advances in topological dynamics, B. - [u. a.], 1973, p. 95-115; [5] M a r c u s B., "Israel J. Math.", 1975, v. 21, № 2-3, p. 133-44; е г 0 \% e, «Invent. math.", 1978, v. 46, №3, p. 201-09; [8] G r ee n L. W., "Duke math. J.", 1974, v. 41, N. 1, p. 115-26; [9] B o w e n R., "Israel J. Math.", 1976, v. 23, № 3-4, p. 26773; [10] B o w e n R., M a r c u s B., там же, 1977, v. 26 , №1, p. 43-67; [11] R a t'ner M., "Ann. Math.», 1982, v. 115, №3,

p. 597-614. ОРЛИЧА КЛАСС - множество функций $L_{M}$, удовлетворяющее условию



\[

\int_{G} M(x(t)) d t<\infty

\]



где $G$ - ограниченное замкнутое множество в $\mathbb{R}^{n}$, $d t$ - мера Лебега, $M(u)$ - четная выпуклая функция, возрастающая при положительных $u$, и



\[

\lim _{u \rightarrow 0} u^{-1} M(u)=\lim _{u \rightarrow \infty} u[M(u)]^{-1}=0 .

\]



Такие функции наз. $N$-ф у н к ц и я м и. Функция $M(u)$ допускает представление



\[

M(u)=\int_{0}^{|u|} p(v) d v

\]



где $p(v)=M^{\prime}(v)$ не убывает на полуоси,



\[

p(0)=\lim _{v \rightarrow 0} p(v)=0,

\]



$p(0)>0$ при $v>0$. Функции $M(u)$ и



\[

N(u)=\int_{0}^{|u|} p^{-1}(v) d v,

\]



где $p^{-1}(v)$ - обратная к $p(v)$ функция, наз. д о п о Јни т ль вы ми фу кция м и. Напр., если





\end{document}